\section{Introduction}

The rapid growth of decentralized finance (DeFi) has led to an explosion of cryptocurrency transactions, with the total value locked in DeFi protocols exceeding \$100 billion as of 2024. This growth has been accompanied by a corresponding increase in fraud, with cross-chain bridge exploits alone causing over \$2 billion in losses between 2021 and 2023~\cite{chainalysis2024}. Detecting illicit transactions in blockchain networks is therefore a critical challenge for maintaining trust in the cryptocurrency ecosystem.

Traditional fraud detection methods rely on hand-crafted features and tabular classifiers~\cite{weber2019anti}, ignoring the rich graph structure of blockchain transaction networks. Recent work has applied Graph Neural Networks (GNNs) to this problem~\cite{pareja2020evolvegcn, weber2019anti}, but existing approaches face two key limitations:

\begin{enumerate}
    \item \textbf{Temporal isolation}: Real-world blockchain datasets such as Elliptic~\cite{weber2019anti} organize transactions into temporal snapshots (timesteps), where edges exist only within the same timestep. Standard GNNs cannot propagate information across timesteps, limiting their ability to capture evolving fraud patterns.

    \item \textbf{Homogeneous edge treatment}: All edges are treated identically, despite representing fundamentally different types of relationships (\eg, direct payments vs.\ cross-temporal similarity patterns).
\end{enumerate}

We propose \textbf{ChainGuard}, a framework built on the Temporal Heterogeneous Graph Neural Network (\ours{}), which addresses both limitations. Our approach makes the following contributions:

\begin{itemize}
    \item \textbf{Temporal $k$-NN graph augmentation}: We introduce cross-timestep edges based on feature-space $k$-nearest-neighbor similarity between adjacent timesteps, breaking the temporal isolation of the original graph. This creates a heterogeneous graph with two edge types: original payment edges and temporal similarity edges.

    \item \textbf{Heterogeneous message passing via R-GCN}: We employ Relational Graph Convolutional Networks~\cite{schlichtkrull2018modeling} to learn separate weight matrices for each edge type, enabling the model to distinguish between direct payment relationships and temporal similarity patterns.

    \item \textbf{Comprehensive ablation study}: Through systematic ablation experiments (M1--M5), we demonstrate the individual contribution of each component and reveal a key insight: \emph{graph augmentation strategy is more important than model complexity}.

    \item \textbf{Extensive baseline comparison}: We evaluate against six baselines spanning non-graph ML (Logistic Regression, Random Forest, Gradient Boosting), standard GNNs (GAT, GraphSAGE), and temporal GNNs (EvolveGCN-H), all under strict temporal evaluation to prevent data leakage.
\end{itemize}

Our best model (M3: R-GCN with temporal $k$-NN augmentation) achieves AUC-ROC of 0.8678 on the Elliptic dataset under temporal split, outperforming all baselines including strong non-graph methods like Random Forest (0.8601) and GraphSAGE (0.8624).
